% Part: lambda-calculus
% Chapter: lambda-definability
% Section: truth-values

\documentclass[../../../include/open-logic-section]{subfiles}

\begin{document}

\olfileid[tr]{lam}{ldf}{tvr}
\olsection{Doğruluk Değerleri ve Bağıntılar}

Doğruluk değerlerini saf lambda hesabında şöyle kodlayabiliriz:
\begin{align*}
  \fn{true} & \ident \lambd[x][\lambd[y][x]]\\
  \fn{false} & \ident \lambd[x][\lambd[y][y]].
\end{align*}

Doğruluk değerleri, iki argüman alıp bunlardan birini döndüren fonksiyonlar,
yani \emph{seçiciler} olarak gösterilir. $\fn{true}$ birinci, $\fn{false}$
ikinci argümanını seçer. Örneğin $\fn{true}\,MN$ her zaman $M$'ye,
$\fn{false}\,MN$ ise her zaman $N$'ye indirgenir.

\begin{defn}
Bir $R\subseteq\Nat^n$ bağıntısına, aşağıdakileri sağlayan bir $R$ terimi
varsa !!{lambda definable} deriz:
\begin{align*}
  R\, \num{n_1} \dots \num{n_n} & \bred \fn{true}
  \intertext{$R(n_1,\dots,n_n)$ geçerli olduğunda ve}
  R\, \num{n_1} \dots \num{n_n} & \bred \fn{false}
\end{align*}
aksi durumda.
\end{defn}

Örneğin yalnızca $0$ için geçerli olan $\fn{IsZero}=\{0\}$ bağıntısı şu
terimle !!{lambda definable}:
\[
  \fn{IsZero} \ident \lambd[n][n (\lambd[x][\fn{false}])\, \fn{true}].
\]
Bu nasıl çalışır? Church sayıları yineleyiciler, yani birinci argümanını
ikinciye $n$ kez uygulayan fonksiyonlar olarak tanımlanır. Başlangıç değerini
$\fn{true}$ seçeriz; yinelemenin her adımında, önceki adımın sonucuna
bakmaksızın $\fn{false}$ döndürürüz. Bu adım başlangıç değerine $n$ kez
uygulanır. Sonuç ancak ve ancak adım hiç uygulanmamışsa, yani $n=0$ ise
$\fn{true}$ olur.

Doğruluk değerlerinin bu gösterimine dayanarak doğruluk fonksiyonlarını da
tanımlayabiliriz. Değilleme ve tümel evetlemenin gösterimleri şöyledir:
\begin{align*}
  \fn{Not} & \ident \lambd[x][x\, \fn{false}\, \fn{true}]\\
  \fn{And} & \ident \lambd[x][\lambd[y][xy \,\fn{false}]].
\end{align*}
$\fn{Not}$ bir argüman alır; argüman $\fn{false}$ ise $\fn{true}$,
$\fn{true}$ ise $\fn{false}$ döndürür. $\fn{And}$ iki doğruluk değeri alıp
ancak ve ancak ikisi de $\fn{true}$ ise $\fn{true}$ döndürmelidir. Doğruluk
değerleri seçici olarak gösterildiğinden, $x$ bir doğruluk değeri olup iki
argümana uygulandığında $x=\fn{true}$ ise birinci, aksi durumda ikinci
argüman sonuç olur. $\fn{And}$, $x,y$ argümanlarını alır ve $y$ ile
$\fn{false}$ değerlerini birinci argümanı $x$'e verir. $x$ bir doğruluk
değeriyse sonuç $x=\fn{true}$ olduğunda $y$, $x=\fn{false}$ olduğunda
$\fn{false}$ olur; istediğimiz de tam olarak budur.

Burada $\fn{And}$ fonksiyonuna yalnızca doğruluk değerlerinin argüman olarak
verildiğini varsayıyoruz. Başka terimler verilirse sonuç, yani varsa normal
biçim, bir doğruluk değeri olmayabilir.

\begin{prob}
  Doğruluk değerlerinin $\lambd$-terimleri olarak kodlanmasını kullanarak
  kapsayıcı ve dışlayıcı veya doğruluk fonksiyonlarını gösteren $\fn{Or}$ ve
  $\fn{Xor}$ fonksiyonlarını tanımlayın.
\end{prob}

\end{document}
